\documentclass[11pt]{article}

\title{Math 603: Homework 3}
\author{Swayam Chube}
\date{Last Updated: \today}

\usepackage[utf8]{inputenc} % allow utf-8 input
\usepackage[T1]{fontenc}    % use 8-bit T1 fonts
\usepackage{hyperref}       % hyperlinks
\usepackage{url}            % simple URL typesetting
\usepackage{booktabs}       % professional-quality tables
\usepackage{amsfonts}       % blackboard math symbols
\usepackage{nicefrac}       % compact symbols for 1/2, etc.
\usepackage{microtype}      % microtypography
\usepackage{graphicx}
\usepackage{natbib}
\usepackage{doi}
\usepackage{amssymb}
\usepackage{bbm}
\usepackage{amsthm}
\usepackage{amsmath}
\usepackage{xcolor}
\usepackage{theoremref}
\usepackage{enumitem}
\usepackage{fouriernc}
\usepackage{mdframed}
\usepackage{mathrsfs}
\setlength{\marginparwidth}{2cm}
\usepackage{todonotes}
\usepackage{stmaryrd}
\usepackage[all,cmtip]{xy} % For diagrams, praise the Freyd-Mitchell theorem 
\usepackage{marvosym}
\usepackage{geometry}
\usepackage{titlesec}
\usepackage{mathtools}
\usepackage{tikz}
\usetikzlibrary{cd}
\usepackage{epigraph}
\setlength\epigraphwidth{0.4\textwidth}

\renewcommand{\qedsymbol}{$\blacksquare$}
% \renewcommand{\familydefault}{\sfdefault} % Do you want this font? 

% Uncomment to override  the `A preprint' in the header
% \renewcommand{\headeright}{}
% \renewcommand{\undertitle}{}
% \renewcommand{\shorttitle}{}

\hypersetup{
    pdfauthor={Swayam Chube},
    colorlinks=true,
	citecolor=blue,
}

\newtheoremstyle{thmstyle}%               % Name
  {}%                                     % Space above
  {}%                                     % Space below
  {}%                             % Body font
  {}%                                     % Indent amount
  {\bfseries\scshape}%                            % Theorem head font
  {.}%                                    % Punctuation after theorem head
  { }%                                    % Space after theorem head, ' ', or \newline
  {\thmname{#1}\thmnumber{ #2}\thmnote{ (#3)}}%                                     % Theorem head spec (can be left empty, meaning `normal')

\newtheoremstyle{defstyle}%               % Name
  {}%                                     % Space above
  {}%                                     % Space below
  {}%                                     % Body font
  {}%                                     % Indent amount
  {\bfseries\scshape}%                            % Theorem head font
  {.}%                                    % Punctuation after theorem head
  { }%                                    % Space after theorem head, ' ', or \newline
  {\thmname{#1}\thmnumber{ #2}\thmnote{ (#3)}}%                                     % Theorem head spec (can be left empty, meaning `normal')

\theoremstyle{thmstyle}
\newtheorem{theorem}{Theorem}[section]
\newtheorem{lemma}[theorem]{Lemma}
\newtheorem{proposition}[theorem]{Proposition}

\theoremstyle{defstyle}
\newtheorem{definition}[theorem]{Definition}
\newtheorem{corollary}[theorem]{Corollary}
\newtheorem{porism}[theorem]{Porism}
\newtheorem{remark}[theorem]{Remark}
\newtheorem{interlude}[theorem]{Interlude}
\newtheorem{example}[theorem]{Example}
\newtheorem*{notation}{Notation}
\newtheorem*{claim}{Claim}

% Common Algebraic Structures
\newcommand{\R}{\mathbb{R}}
\newcommand{\Q}{\mathbb{Q}}
\newcommand{\Z}{\mathbb{Z}}
\newcommand{\N}{\mathbb{N}}
\newcommand{\bbC}{\mathbb{C}} 
\newcommand{\K}{\mathbb{K}} % Base field which is either \R or \bbC
\newcommand{\F}{\mathbb{F}} % Base field which is either \R or \bbC
\newcommand{\calA}{\mathcal{A}} % Banach Algebras
\newcommand{\calB}{\mathcal{B}} % Banach Algebras
\newcommand{\calI}{\mathcal{I}} % ideal in a Banach algebra
\newcommand{\calJ}{\mathcal{J}} % ideal in a Banach algebra
\newcommand{\calR}{\mathcal{R}} % range of an operator
\newcommand{\frakM}{\mathfrak{M}} % sigma-algebra
\newcommand{\calO}{\mathcal{O}} % Ring of integers
\newcommand{\bbA}{\mathbb{A}} % Adele (or ring thereof)
\newcommand{\bbI}{\mathbb{I}} % Idele (or group thereof)

% Categories
\newcommand{\catTopp}{\mathbf{Top}_*}
\newcommand{\catGrp}{\mathbf{Grp}}
\newcommand{\catTopGrp}{\mathbf{TopGrp}}
\newcommand{\catSet}{\mathbf{Set}}
\newcommand{\catTop}{\mathbf{Top}}
\newcommand{\catRing}{\mathbf{Ring}}
\newcommand{\catCRing}{\mathbf{CRing}} % comm. rings
\newcommand{\catMod}{\mathbf{Mod}}
\newcommand{\catMon}{\mathbf{Mon}}
\newcommand{\catMan}{\mathbf{Man}} % manifolds
\newcommand{\catDiff}{\mathbf{Diff}} % smooth manifolds
\newcommand{\catAlg}{\mathbf{Alg}}
\newcommand{\catRep}{\mathbf{Rep}} % representations 
\newcommand{\catVec}{\mathbf{Vec}}

% Group and Representation Theory
\newcommand{\chr}{\operatorname{char}}
\newcommand{\Aut}{\operatorname{Aut}}
\newcommand{\GL}{\operatorname{GL}}
\newcommand{\im}{\operatorname{im}}
\newcommand{\tr}{\operatorname{tr}}
\newcommand{\id}{\mathbf{id}}
\newcommand{\cl}{\mathbf{cl}}
\newcommand{\Gal}{\operatorname{Gal}}
\newcommand{\Tr}{\operatorname{Tr}}
\newcommand{\sgn}{\operatorname{sgn}}
\newcommand{\Sym}{\operatorname{Sym}}
\newcommand{\Alt}{\operatorname{Alt}}

% Commutative and Homological Algebra
\newcommand{\spec}{\operatorname{spec}}
\newcommand{\mspec}{\operatorname{m-spec}}
\newcommand{\Spec}{\operatorname{Spec}}
\newcommand{\MaxSpec}{\operatorname{MaxSpec}}
\newcommand{\Tor}{\operatorname{Tor}}
\newcommand{\tor}{\operatorname{tor}}
\newcommand{\Ann}{\operatorname{Ann}}
\newcommand{\Supp}{\operatorname{Supp}}
\newcommand{\Hom}{\operatorname{Hom}}
\newcommand{\End}{\operatorname{End}}
\newcommand{\coker}{\operatorname{coker}}
\newcommand{\limit}{\varprojlim}
\newcommand{\colimit}{%
  \mathop{\mathpalette\colimit@{\rightarrowfill@\textstyle}}\nmlimits@
}
\makeatother


\newcommand{\fraka}{\mathfrak{a}} % ideal
\newcommand{\frakb}{\mathfrak{b}} % ideal
\newcommand{\frakc}{\mathfrak{c}} % ideal
\newcommand{\frakf}{\mathfrak{f}} % face map
\newcommand{\frakg}{\mathfrak{g}}
\newcommand{\frakh}{\mathfrak{h}}
\newcommand{\frakm}{\mathfrak{m}} % maximal ideal
\newcommand{\frakn}{\mathfrak{n}} % naximal ideal
\newcommand{\frakp}{\mathfrak{p}} % prime ideal
\newcommand{\frakq}{\mathfrak{q}} % qrime ideal
\newcommand{\fraks}{\mathfrak{s}}
\newcommand{\frakt}{\mathfrak{t}}
\newcommand{\frakz}{\mathfrak{z}}
\newcommand{\frakA}{\mathfrak{A}}
\newcommand{\frakB}{\mathfrak{B}}
\newcommand{\frakI}{\mathfrak{I}}
\newcommand{\frakJ}{\mathfrak{J}}
\newcommand{\frakK}{\mathfrak{K}}
\newcommand{\frakL}{\mathfrak{L}}
\newcommand{\frakN}{\mathfrak{N}} % nilradical 
\newcommand{\frakO}{\mathfrak{O}} % dedekind domain
\newcommand{\frakP}{\mathfrak{P}} % Prime ideal above
\newcommand{\frakQ}{\mathfrak{Q}} % Qrime ideal above 
\newcommand{\frakR}{\mathfrak{R}} % jacobson radical
\newcommand{\frakU}{\mathfrak{U}}
\newcommand{\frakV}{\mathfrak{V}}
\newcommand{\frakW}{\mathfrak{W}}
\newcommand{\frakX}{\mathfrak{X}}

% General/Differential/Algebraic Topology 
\newcommand{\scrA}{\mathscr{A}}
\newcommand{\scrB}{\mathscr{B}}
\newcommand{\scrF}{\mathscr{F}}
\newcommand{\scrM}{\mathscr{M}}
\newcommand{\scrN}{\mathscr{N}}
\newcommand{\scrP}{\mathscr{P}}
\newcommand{\scrO}{\mathscr{O}} % sheaf
\newcommand{\scrR}{\mathscr{R}}
\newcommand{\scrS}{\mathscr{S}}
\newcommand{\bbH}{\mathbb H}
\newcommand{\Int}{\operatorname{Int}}
\newcommand{\psimeq}{\simeq_p}
\newcommand{\wt}[1]{\widetilde{#1}}
\newcommand{\RP}{\mathbb{R}\text{P}}
\newcommand{\CP}{\mathbb{C}\text{P}}

% Miscellaneous
\newcommand{\wh}[1]{\widehat{#1}}
\newcommand{\calM}{\mathcal{M}}
\newcommand{\calP}{\mathcal{P}}
\newcommand{\onto}{\twoheadrightarrow}
\newcommand{\into}{\hookrightarrow}
\newcommand{\Gr}{\operatorname{Gr}}
\newcommand{\Span}{\operatorname{Span}}
\newcommand{\ev}{\operatorname{ev}}
% \newcommand{\weakto}{\stackrel{w}{\longrightarrow}}
\newcommand{\weakto}{\rightharpoonup}

\newcommand{\define}[1]{\textcolor{blue}{\textit{#1}}}
% \newcommand{\caution}[1]{\textcolor{red}{\textit{#1}}}
\newcommand{\important}[1]{\textcolor{red}{\textit{#1}}}
\renewcommand{\mod}{~\mathrm{mod}~}
\renewcommand{\le}{\leqslant}
\renewcommand{\leq}{\leqslant}
\renewcommand{\ge}{\geqslant}
\renewcommand{\geq}{\geqslant}
\newcommand{\Res}{\operatorname{Res}}
\newcommand{\floor}[1]{\left\lfloor #1\right\rfloor}
\newcommand{\ceil}[1]{\left\lceil #1\right\rceil}
\newcommand{\gl}{\mathfrak{gl}}
\newcommand{\ad}{\operatorname{ad}}
\newcommand{\Stab}{\operatorname{Stab}}
\newcommand{\bfX}{\mathbf{X}}
\newcommand{\Ind}{\operatorname{Ind}}
\newcommand{\bfG}{\mathbf{G}}
\newcommand{\rank}{\operatorname{rank}}
\newcommand{\calo}{\mathcal{o}}
\newcommand{\frako}{\mathfrak{o}}
\newcommand{\Cl}{\operatorname{Cl}}

\newcommand{\idim}{\operatorname{idim}}
\newcommand{\pdim}{\operatorname{pdim}}
\newcommand{\Ext}{\operatorname{Ext}}
\newcommand{\co}{\operatorname{co}}
\newcommand{\bfO}{\mathbf{O}}
\newcommand{\bfF}{\mathbf{F}} % Fitting Subgroup
\newcommand{\Syl}{\operatorname{Syl}}
\newcommand{\nor}{\vartriangleleft}
\newcommand{\noreq}{\trianglelefteqslant}
\newcommand{\subnor}{\nor\!\nor}
\newcommand{\Soc}{\operatorname{Soc}}
\newcommand{\core}{\operatorname{core}}
\newcommand{\Sd}{\operatorname{Sd}}
\newcommand{\mesh}{\operatorname{mesh}}
\newcommand{\sminus}{\setminus}
\newcommand{\diam}{\operatorname{diam}}
\newcommand{\Ass}{\operatorname{Ass}}
\newcommand{\projdim}{\operatorname{proj~dim}}
\newcommand{\injdim}{\operatorname{inj~dim}}
\newcommand{\gldim}{\operatorname{gl~dim}}
\newcommand{\embdim}{\operatorname{emb~dim}}
\newcommand{\hght}{\operatorname{ht}}
\newcommand{\depth}{\operatorname{depth}}
\newcommand{\ul}[1]{\underline{#1}}
\newcommand{\type}{\operatorname{type}}
\newcommand{\CM}{\operatorname{CM}}
\newcommand{\cech}[1]{\mathbin{\check{#1}}}
\newcommand{\cdim}{\operatorname{cdim}}
\newcommand{\Der}{\operatorname{Der}}
\newcommand{\trdeg}{\operatorname{trdeg}}
\newcommand{\calE}{\mathcal{E}}
\newcommand{\calF}{\mathcal{F}}
\newcommand{\calT}{\mathcal{T}}
\newcommand{\calN}{\mathcal{N}}
\newcommand{\calL}{\mathcal{L}}
\renewcommand{\Re}{\operatorname{Re}}
\renewcommand{\Im}{\operatorname{Im}}
\newcommand{\gr}{\operatorname{gr}}
\newcommand{\dist}{\operatorname{dist}}
\newcommand{\calH}{\mathcal{H}}


\geometry {
    margin = 1in
}

\titleformat
{\section}
[block]
{\Large\bfseries\sffamily}
{\S\thesection}
{0.5em}
{\centering}
[]


\titleformat
{\subsection}
[block]
{\normalfont\bfseries\sffamily}
{\S\S}
{0.5em}
{\centering}
[]

\begin{document}
\maketitle 

\section{Problem 1}

\subsection{Part (a)}

That $T$ is linear is clear. We must show that it is an isometric bijection. We have 
\begin{equation*}
    |Ty(x)| = \left|\sum_{n = 1}^\infty x_ny_n\right|\le\sum_{n = 1}^\infty |x_ny_n|\le \|y\|_\infty\sum_{n = 1}^\infty |x_n| = \|x\|_1\|y\|_\infty,
\end{equation*}
so that $\|Ty\|\le \|y\|_\infty$ for all $y\in\ell^\infty$. In particular, $T$ is a bounded linear functional with $\|T\|\le 1$.

For each $n\ge 1$, let $e_n$ denote the sequence 
\begin{equation*}
    e_n(m) = \begin{cases}
        1 & n = m\\
        0 & \text{otherwise}.
    \end{cases}
\end{equation*}
If $y = (y_n)_{n = 1}^\infty\in\ell^\infty$ is such that $Ty = 0$, then $y_n = Ty(e_n) = 0$ for all $n\ge 1$, so that $y = 0$. This shows that $T$ has trivial kernel, that is, $T$ is injective.

Next, we show that $T$ is bijective. Let $\varphi\in\left(\ell^1\right)^\ast$ and set $a_n = \varphi(e_n)$ for all $n\ge 1$. Then 
\begin{equation*}
    |a_n| = |\varphi(e_n)|\le\|\varphi\|\|e_n\|_1 = \|\varphi\|,
\end{equation*}
so that $a = (a_n)_{n = 1}^\infty\in\ell^\infty$. We shall show that $\varphi = Ta$. Indeed, for any $x = (x_n)_{n = 1}^\infty\in\ell^1$ and $N\ge 1$, define $x^N\in\ell^1$ as 
\begin{equation*}
    x^N_n = \begin{cases}
        x_n & n\le N\\
        0 & n > N.
    \end{cases}
\end{equation*}
Then 
\begin{equation*}
    \|x - x^N\|\le\sum_{n = N + 1}^\infty |x_n|,
\end{equation*}
so that $\|x - x^N\|\to 0$ as $N\to\infty$, since the tail sum of a convergent sequence goes to $0$. In other words, $x^N\to x$ as $N\to\infty$. Further, note that 
\begin{equation*}
    (Ta)(x^N) = \sum_{n = 1}^N a_n x_n = \sum_{n = 1}^N \varphi(x_n e_n) = \varphi\left(\sum_{n = 1}^N x_n e_n\right) = \varphi(x^N).
\end{equation*}
Since $\varphi$ and $Ta$ are continuous functions, we have 
\begin{equation*}
    \varphi(x) = \lim_{N\to\infty} \varphi(x^N) = \lim_{N\to\infty} (Ta)(x^N) = (Ta)(x),
\end{equation*}
so that $\varphi = Ta$, as desired. This shows that $T$ is a continuous linear bijection between Banach spaces, and hence, due to the Open Mapping Theorem (or the Bounded Inverse Theorem), $T$ is an isomorphism of Banach spaces.

It remains to show that $T$ is an isometry, to which end, we must show that $\|Ty\| = \|y\|_\infty$ for all $y\in\ell^\infty$. Note that we have already shown that $\|Ty\|\le\|y\|_\infty$. Further, for $y = (y_n)_{n = 1}^\infty\in\ell^\infty$, note that $(Ty)(e_n) = y_n$ for all $n\ge 1$. Since $\|e_n\|_1 = 1$, we have that $\|Ty\|\ge |y_n|$ for all $n\ge 1$, consequently, $\|Ty\|\ge \|y\|_\infty$. This shows that $\|Ty\| = \|y\|_\infty$, thereby completing the proof.

\subsection{Part (b)}

That $S$ is linear is clear. We must show that it is an isometric bijection. We have 
\begin{equation*}
    |(Sy)(x)| = \left|\sum_{n = 1}^\infty x_ny_n\right|\le\sum_{n = 1}^\infty |x_ny_n|\le \|x\|_\infty\sum_{n = 1}^\infty |y_n| = \|x\|_\infty \|y\|_1,
\end{equation*}
so that $\|Sy\|\le \|y\|_1$. As in the preceding part, consider the ``standard basis vectors'' $e_n$ for $n\ge 1$. If $y\in\ell^1$ is such that $Sy = 0$, then for all $n\ge 1$, we have $y_n = (Sy)(e_n) = 0$, so that $y = 0$. This shows that $S$ is injective. 

Next, we show that $S$ is surjective. Let $\varphi\in (c_0)^\ast$. Set $a_n = \varphi(e_n)$, and $a = (a_n)_{n =1 }^\infty$. First, we shall show that $a\in\ell^1$. To this end, define 
\begin{equation*}
    \zeta_n = 
    \begin{cases}
        \frac{\overline a_n}{|a_n|} & a_n\ne 0\\
        0 & a_n = 0.
    \end{cases}
\end{equation*}
Note that $\zeta_n a_n = |a_n|$ and $|\zeta_n|\le 1$. Further, for every positive integer $N$, $\displaystyle\sum_{n = 1}^N \zeta_n e_n\in c_0$, so that 
\begin{equation*}
    \sum_{n = 1}^N|a_n| = \sum_{n = 1}^N \zeta_n a_n = \sum_{n = 1}^N \zeta_n\varphi(e_n) = \varphi\left(\sum_{n = 1}^N \zeta_n e_n\right)\le \|\varphi\|,
\end{equation*}
since $\displaystyle\left\|\sum_{n = 1}^N \zeta_n e_n\right\|\le 1$. Taking the limit $N\to\infty$, we have 
\begin{equation*}
    \|a\|_1 = \sum_{n = 1}^\infty |a_n|\le\|\varphi\|.
\end{equation*}
We contend that $\varphi = Sa$. Indeed, for any $x = (x_n)_{n = 1}^\infty\in c_0$ and $N\ge 1$, define $x^N\in c_0$ as 
\begin{equation*}
    x^N_n = \begin{cases}
        x_n & n\le N\\
        0 & n > N.
    \end{cases}
\end{equation*}
We have 
\begin{equation*}
    \|x - x^N\| = \sup_{n\ge N + 1} |x_n|.
\end{equation*}
But since $x\in c_0$, by definition, we have that the right hand side above goes to $0$ as $N\to\infty$. This means $x^N\to x$ in $c_0$.

Clearly, we have 
\begin{equation*}
    (Sa)(x^N) = \sum_{n = 1}^N a_n x_n = \sum_{n = 1}^N x_n\varphi(e_n) = \varphi\left(\sum_{n = 1}^N x_n e_n\right) = \varphi(x^N).
\end{equation*}
Since $Sa$ and $\varphi$ are continuous, taking the limit of the above equality as $N\to\infty$ and using that $x^N\to x$, we have 
\begin{equation*}
    (Sa)(x) = \lim_{N\to\infty} (Sa)(x^N) = \lim_{N\to\infty}\varphi(x^N) = \varphi(x),
\end{equation*}
and hence $\varphi = Sa$, as desired. 

Thus $S$ is a bijective bounded linear map between Banach spaces, and due to the Open Mapping Theorem (or the Bounded Inverse Theorem), $S$ must be an isomorphism.

It remains to show that $S$ is an isometry. We have already shown that $\|Sy\|\le \|y\|_1$ for all $y\in\ell^1$. To see that the reverse inequality is also true, for $y = (y_n)_{n = 1}^\infty\in\ell^1$, define 
\begin{equation*}
    \zeta_n = 
    \begin{cases}
        \frac{y_n}{|y_n|} & y_n\ne 0\\
        0 & y_n = 0.
    \end{cases}
\end{equation*}
Note that $|\zeta_n|\le 1$ and $\zeta_ny_n = |y_n|$. As a result, 
\begin{equation*}
    \sum_{n = 1}^N |y_n| = \sum_{n = 1}^N \zeta_n y_n = \sum_{n = 1}^N \zeta_n (Sy)(e_n) = (Sy)\left(\sum_{n = 1}^N\zeta_n e_n\right).
\end{equation*}
But using the fact that $\left\|\displaystyle\sum_{n = 1}^N \zeta_n e_n\right\|_\infty\le 1$, we have 
\begin{equation*}
    \sum_{n = 1}^N |y_n|\le \|Sy\|,
\end{equation*}
and hence, in the limit $N\to\infty$, we obtain $\|Sy\|\ge \|y\|_1$, which shows that $\|Sy\| = \|y\|_1$ for all $y\in\ell^1$, thereby completing the proof that $S$ is an isometric bijection. 

\subsection{Part (c)}

The sequence $(1, 1, 1, \dots)$ is clearly an element of $\ell^\infty$ but not an element of $c_0$. Therefore, it only remains to show that $c_0$ is closed in the $\ell^\infty$-norm. Indeed, suppose $(x^N)_{N\ge 1}$ is a convergent sequence in $\ell^\infty$ with $x^N\in c_0$ for all $N\ge 1$ and let $x^N\to x = (x_n)_{n = 1}^\infty$ in $\ell^\infty$. We shall show that $x\in c_0$.

Let $\varepsilon > 0$. Then there is a positive integer $M\ge 1$ such that for all $N\ge M$, $\|x - x^N\|_\infty < \frac{\varepsilon}{2}$. Consider $x^M = (x^M_n)_{n = 1}^\infty$. Since this is an element of $c_0$, there is a positive integer $\wt M\ge 1$ such that for all $n\ge\wt M$, $|x^M_n| < \frac{\varepsilon}{2}$. Now, for $n\ge\wt M$, using the triangle inequality, we can write 
\begin{equation*}
    \frac{\varepsilon}{2} > \|x - x^M\|\ge |x_n - x^M_n|\ge|x_n| - |x^M_n|,
\end{equation*}
so that 
\begin{equation*}
    |x_n| < \frac{\varepsilon}{2} + |x^M_n| < \varepsilon.
\end{equation*}
Thus $x\in c_0$, as desired, and $c_0$ is a proper closed subspace of $\ell^\infty$.

\subsection{Part (d)}

We have seen in part (c) that $c_0\subsetneq \ell^\infty$ is a proper closed subspace. Thus, as we have seen in class, there is a bounded linear functional $0\ne\varphi\in\left(\ell^\infty\right)^\ast$ such that $\varphi|_{c_0} = 0$. Suppose now that $\ell^1$ is reflexive, that is, the canonical evaluation map $J\colon\ell^1\to\left(\ell^1\right)^{\ast\ast}$ is an isometric isomorphism (to which end, it suffices for $J$ to be surjective since it is always an isometry). 
Note that through the isometric isomorphism $T\colon\ell^\infty\to\left(\ell^1\right)^\ast$, the bounded linear map $\varphi\circ T^{-1}$ is a bounded linear functional on $\left(\ell^1\right)^\ast$. 

Since we have assumed $J$ to be surjective, there exists an $x\in\ell^1$ such that $\varphi\circ T^{-1} = J(x)$. Note that for every $y = (y_n)_{n = 1}^\infty\in c_0\subseteq\ell^\infty$, we have 
\begin{equation*}
    \sum_{n = 1}^\infty x_ny_n = (Ty)(x) = J(x)\left(T(y)\right) = \left(\varphi\circ T^{-1}\right)\left(T(y)\right) = \varphi(y) = 0.
\end{equation*}
Since $e_n\in c_0$ for all $n\ge 1$, we have that 
\begin{equation*}
    x_n = (Te_n)(x) = 0
\end{equation*}
for all $n\ge 1$, so that $x = 0$, i.e., $\varphi = J(x) = 0$, a contradiction. Hence, $\ell^1$ is not reflexive, as desired.

\section{Problem 2}

\subsection{Part (a)}

The sequence $e_n$ never converges strongly for any $1\le p < \infty$. Indeed, for $m\ne n$, we have 
\begin{equation*}
    \|e_m - e_n\| = \left(1^p + 1^p\right)^{\frac{1}{p}} = 2^{\frac{1}{p}}.
\end{equation*}
Since all convergent sequences are Cauchy, it follows from the above equality that $(e_n)_{n = 1}^\infty$ cannot be Cauchy, and therefore, cannot be convergent. 

\subsection{Part (b)}

We shall show that $(e_n)_{n = 1}^\infty$ converges weakly for $1 < p < \infty$ but not for $p = 1$.

First, we show that $(e_n)_{n = 1}^\infty$ does not converge weakly in $\ell^1$. Recall from Problem 1 (a) that the map $T\colon \ell^\infty\to\left(\ell^1\right)^\ast$ is an isometric isomorphism of Banach spaces. Let $z\in\ell^\infty$ denote the sequence $z_n = 1$ for all $n\ge 1$. Suppose $e_n\weakto x$ for some $x = (x_n)_{n = 1}^\infty\in\ell^1$. Then, since $e_n\in\ell^\infty$, we must have that for all $n\ge 1$, 
\begin{equation*}
    x_n = (Te_n)(x) = \lim_{m\to\infty} (Te_n)(e_m) = 0,
\end{equation*}
since $(Te_n)(e_m) = 0$ for $n > m$. Hence, $x = 0$. But then we must have that 
\begin{equation*}
    0 = (Tz)(x) = \lim_{m\to\infty} (Tz)(e_m) = 1,
\end{equation*}
a contradiction. Thus $(e_n)_{n = 1}^\infty$ does not converge weakly in $\ell^1$.

Suppose now that $1 < p < \infty$. We shall show that $(e_n)_{n = 1}^\infty$ converges weakly in $\ell^p$. In particular, we shall show that $e_n\weakto 0$. Let $\varphi\in\left(\ell^p\right)^\ast$ be a bounded linear functional, and set $a_n = \varphi(e_n)$. Let $q$ denote the conjugate exponent to $p$, that is, $q = \frac{p}{p - 1}$, and set 
\begin{equation*}
    b_n = \begin{cases}
        |a_n|^{q - 2}\overline a_n & a_n\ne 0\\
        0 & a_n = 0,
    \end{cases}
\end{equation*}
so that $a_nb_n = |a_n|^{q}$ for all $n\ge 1$.
We shall show that $\displaystyle\sum_{n = 1}^\infty |a_n|^{q} < \infty$. Indeed, for any positive integer $N$, we have 
\begin{align*}
    \sum_{n = 1}^N |a_n|^q &= \sum_{n = 1}^N a_nb_n\\ &= \sum_{n = 1}^N b_n\varphi(e_n)\\ &= \varphi\left(\sum_{n = 1}^N b_ne_n\right)\\ &= \left|\varphi\left(\sum_{n = 1}^N b_ne_n\right)\right|\\ &\le\|\varphi\|\left(\sum_{n = 1}^N |b_n|^p\right)^{\frac{1}{p}}\\ &= \|\varphi\|\left(\sum_{n = 1}^N |a_n|^{p(q - 1)}\right)^{\frac{1}{p}}\\
    &= \|\varphi\|\left(\sum_{n = 1}^N |a_n|^q\right),
\end{align*}
so that 
\begin{equation*}
    \left(\sum_{n = 1}^N |a_n|^q\right)^{\frac{1}{p}}\le\|\varphi\|.
\end{equation*}
Taking the limit $N\to\infty$, it follows that $\displaystyle\sum-{n = 1}^\infty |a_n|^q$ converges, so that $a_n\to 0$ as $n\to\infty$. In other words, $e_n\weakto 0$, as desired.

\section{Problem 3}

First note that due to the Parallelogram Law, we have 
\begin{equation*}
    2\left(\left\|\frac{x + y}{2}\right\|^2 + \left\|\frac{x - y}{2}\right\|^2\right) = \|x\|^2 + \|y\|^2.
\end{equation*}
Now let $\|x\|\le 1$ and $\|y\|\le 1$ and $2 \ge \varepsilon > 0$ be such that $\|x - y\|\ge\varepsilon$. Then, we have the inequalities
\begin{equation*}
    2\ge\|x\|^2 + \|y\|^2 = 2\left(\left\|\frac{x + y}{2}\right\|^2 + \left\|\frac{x - y}{2}\right\|^2\right) > 2\left(\left\|\frac{x + y}{2}\right\|^2 + \frac{\varepsilon^2}{4}\right),
\end{equation*}
so that 
\begin{equation*}
    \left\|\frac{x + y}{2}\right\|\le\sqrt{1 - \frac{\varepsilon^2}{4}} < 1 - \frac{\varepsilon^2}{8},
\end{equation*}
where the inequality on the right is evident from squaring both side of the inequality. Thus, taking $\delta = \frac{\varepsilon^2}{\delta}$, we see that the Hilbert space is uniformly convex.

\section{Problem 4}

\begin{lemma}
    Let $\calH$ be a Hilbert space, and $A\in\calL(\calH)$. Then 
    \begin{equation*}
        \|A\| = \sup_{\|x\| = \|y\| = 1}|\langle Ax, y\rangle|.
    \end{equation*}
\end{lemma}
\begin{proof}
    Let $B$ denote the right hand side of the above equality. Due to the Cauchy-Schwarz inequality, for any $\|x\| = \|y\| = 1$, we have 
    \begin{equation*}
        |\langle Ax, y\rangle|\le \|Ax\|\|y\|\le \|A\|\|x\|\|y\| = \|A\|,
    \end{equation*}
    so that $B\le \|A\|$. On the other hand, we have 
    \begin{equation*}
        B\ge\sup_{\substack{\|x\| = 1\\ Ax\ne 0}}\left|\left\langle Ax, \frac{Ax}{\|Ax\|}\right\rangle\right| = \sup_{\substack{\|x\| = 1\\ Ax\ne 0}}\|Ax\| = \sup_{\|x\| = 1}\|Ax\| = \|A\|,
    \end{equation*}
    so that $B = \|A\|$, as desired.
\end{proof}

Coming back to our problem at hand, we have 
\begin{equation*}
    \|A^\ast A\| = \sup_{\|x\| = \|y\| = 1}|\langle A^\ast A x, y\rangle| = \sup_{\|x\| = \|y\| = 1}|\langle Ax, Ay\rangle|\ge\sup_{\|x\| = 1}|\langle Ax, Ax\rangle| = \sup_{\|x\| = 1}\|Ax\|^2 = \|A\|^2.
\end{equation*}
Conversely, again, using the Cauchy-Schwarz Inequality, we have 
\begin{equation*}
    \|A^\ast A\| = \sup_{\|x\| = \|y\| = 1}|\langle A^\ast A x, y\rangle| = \sup_{\|x\| = \|y\| = 1}|\langle Ax, Ay\rangle|\le\sup_{\|x\| = \|y\| = 1}\|Ax\|\|Ay\|\le\sup_{\|x\| = \|y\| = 1}\|A\|^2 = \|A\|^2.
\end{equation*}
Thus $\|A^\ast A\| = \|A\|^2$.

\section{Problem 5}

First, we show that for any projection operator $P$, $\ker P = \calR(I - P)$. Since $P^2 = P$, we have $P(I - P) = 0$, so that $\calR(I - P)\subseteq\ker P$. Conversely, suppose $x\in\ker P$. Due to the Projection Theorem (or corollary thereof), there is a decomposition $\calH = \ker(I - P)\oplus\ker(I - P)^\perp$, so that we can write $x = y + z$ where $y\in\ker (I - P)$ and $z\in\ker(I - P)^\perp$. We can write 
\begin{equation*}
    x = x - Px = (I - P)x = (I - P)y + (I - P)z = (I - P)z,
\end{equation*}
i.e., $x\in\calR(I - P)$, so that $\calR(I - P) = \ker P$. Similarly, since $I - P$ is a projection, we also have $\calR(P) = \ker(I - P)$.

In view of the above discussion, (c) is equivalent to the statement ``$\ker P$ and $\calR(P)$ are orthogonal''. We shall use this equivalent formulation in proving the proposition.

\begin{proof}[Proof of Proposition]
    $\boxed{(a)\implies(c)}$ Suppose $P$ is self-adjoint. Then any element of $\calR(P)$ is of the form $Px$ for some $x\in\calH$ and every element of $\calR(I - P)$ is of the form $(I - P)x$ for some $y\in\calH$. Note that 
    \begin{equation*}
        \langle Px, (I - P)y\rangle = \langle(I - P^\ast)Px, y\rangle = \langle (I - P)Px, y\rangle = 0,
    \end{equation*}
    since $P - P^2 = 0$.

    $\boxed{(c)\implies(a)}$ Since $\calR(P)$ and $\calR(I - P)$ are orthogonal, we have that for every $x, y\in\calH$, 
    \begin{equation*}
        \langle Px, (I - P)y\rangle = 0.
    \end{equation*}
    We may rewrite the above as 
    \begin{equation*}
        \langle (I - P^\ast)Px, y\rangle = 0\qquad\text{for all }x, y\in\calH.
    \end{equation*}
    In particular, taking $y = (I - P^\ast)Px$, we see that $(I - P^\ast)Px = 0$ for all $x\in\calH$, that is, $(I - P^\ast)P = 0$, equivalently, $P = P^\ast P$.

    On the other hand, we can write 
    \begin{equation*}
        0 = \overline{\langle Px, (I - P)y\rangle} = \langle (I - P)y, Px\rangle = \langle P^\ast(I - P)y, x\rangle\qquad\text{for all } x, y\in\calH.
    \end{equation*}
    In particular, taking $x = P^\ast(I - P)y$, we see that $P^\ast(I - P)y = 0$ for all $y\in\calH$, that is, $P^\ast(I - P) = 0$, equivalently, $P^\ast = P^\ast P$. This shows that $P = P^\ast$, i.e., $P$ is self-adjoint.

    $\boxed{(a)\implies(b)}$ Suppose $P$ is self-adjoint. Using Problem 4, we have 
    \begin{equation*}
        \|P\| = \|P^2\| = \|P^\ast P\| = \|P\|^2,
    \end{equation*}
    so $\|P\| = 0$ or $\|P\| = 1$. But since $P$ is a non-zero projection operator, it follows that $\|P\| = 1$. 

    $\boxed{(b)\implies(c)}$ We shall use the equivalent formulation of (c) in the discussion preceding the proof. As such, it suffices to show that $\calR(P)\perp\ker P$. Suppose $\|P\| = 1$, and let $x\in\calR(P)$ and $y\in\ker P$. Note that since $P^2 = P$, we have $P(x) = x$. Thus for any $\lambda\in\K$, we have 
    \begin{equation*}
        \|x\|^2 = \|P(x)\|^2 = \|P(x + \lambda y)\|^2\le \|x + \lambda y\|^2 = \|x\|^2 + |\lambda|^2\|y\|^2 + \langle x, \lambda y\rangle + \langle \lambda y, x\rangle,
    \end{equation*}
    therefore, 
    \begin{equation*}
        0\le |\lambda|^2\|y\|^2 + \langle x, \lambda y\rangle + \langle \lambda y, x\rangle = |\lambda|^2\|y\|^2 + 2\Re\left(\langle x, \lambda y\rangle\right) = |\lambda|^2\|y\|^2 + 2\Re\left(\overline\lambda\langle x, y\rangle\right).
    \end{equation*}
    If $y = 0$, then $\langle x, y\rangle = 0$ trivially, and there is nothing to prove. Suppose $y\ne 0$. Setting $\lambda = -\frac{\overline{\langle x, y\rangle}}{\|y\|^2}$, we obtain 
    \begin{equation*}
        0\le - \frac{|\langle x, y\rangle|^2}{\|y\|^2},
    \end{equation*}
    which is possible if and only if $\langle x, y\rangle = 0$. Thus $\calR(P)\perp\ker P$, thereby completing the proof.
\end{proof}

\section{Problem 6}

We recall from class that $x_n\to x$ in a topology induced by functions $\{f_\alpha\colon X\to Y_\alpha\}_{\alpha\in\scrA}$ on $X$ if and only if $f_\alpha(x_n)\to f_\alpha(x)$ In $Y_\alpha$ for each $\alpha\in\scrA$. We shall use this observation repeatedly in the following analysis.

\subsection{Part (a)}

Note that convergence in the strong operator topology is the same as pointwise convergence of these bounded linear operators. Suppose $A_n\to A$ in the operator topology, that is, $\|A - A_n\|\to 0$ as $n\to\infty$. Then for any $x\in X$, we have the inequality: 
\begin{equation*}
    \|Ax - A_n x\| = \|(A - A_n)x\|\le \|A - A_n\|\|x\|.
\end{equation*}
Since the right hand side of the above inequality goes to $0$ as $n\to\infty$, we have that $\|Ax - A_n x\|\to 0$ as $n\to\infty$, i.e., $A_n x\to Ax$ as $n\to\infty$. In other words, $A_n\to A$ in the strong operator topology of $\calL(X)$.

The converse is not true. Let $\tau\colon\ell^2\to\ell^2$ denote the ``left-shift operator'', that is, 
\begin{equation*}
    \tau(x_1, x_2, \dots) = \left(x_2, x_3,  \dots\right).
\end{equation*}
This is clearly a bounded linear operator with $\|\tau\| = 1$. Define $A_n = \tau^{\circ n}$ for each positive integer $n\ge 1$. Note that 
\begin{equation*}
    A_n\left(x_1, x_2, \dots\right) = \left(x_{n + 1}, x_{n + 2},\dots\right).
\end{equation*}
We note that $\|A_n\| = 1$ for every $n\ge 1$. Indeed, since $\displaystyle\sum_{k = n + 1}^\infty |x_k|^2\le\sum_{k = 1}^\infty |x_k|^2$, we have that $\|A_n\|\le 1$, and since 
\begin{equation*}
    A_n(\underbrace{0, \dots, 0}_{n\text{ times}}, 1, 0, 0, \dots) = \left(1, 0, 0, \dots\right),
\end{equation*}
so that $\|A_n\| = 1$. On the other hand, for any $x = (x_n)_{n = 1}^\infty\in\ell^2$, we claim that $A_n x\to 0$. Indeed, 
\begin{equation*}
    \|A_n x\|^2 = \sum_{k = n + 1}^\infty |x_k|^2,
\end{equation*}
which is the tail-sum of an (absolutely) convergent infinite sum, and therefore, it must go to $0$ as $n\to\infty$. Hence, $A_n\to 0$ in the strong operator topology. But since $\|A_n\| = 1$ for all $n\ge 1$, we cannot have $A_n\to 0$ in the operator topology.

\subsection{Part (b)}

Note that $A_n\to A$ in the weak operator topology if and only if for every $x\in X$ and $\varphi\in X^\ast$, $\varphi(A_n x)\to \varphi(Ax)$. Since $A_n\to A$ in the strong operator topology, we have that $A_nx\to Ax$ in $X$. Further since $\varphi$ is a continuous function, we must have that $\varphi(A_n x)\to \varphi(Ax)$ in $\K$, as desired. Thus $A_n\to A$ in the weak operator topology.

The converse is not true, let $\sigma\colon \ell^2\to\ell^2$ denote the ``right-shift operator'', that is, 
\begin{equation*}
\sigma(x_1, x_2, \dots) = (0, x_1, x_2, \dots).
\end{equation*}
Note that $\sigma$ is a bounded-linear operator with $\|\sigma\| = 1$. Define $A_n = \sigma^{\circ n}$ for every positive integer $n\ge 1$. We have 
\begin{equation*}
    A_n(x_1, x_2,\dots) = (\underbrace{0, \dots, 0}_{n\text{ times}}, x_1, x_2, \dots).
\end{equation*}
Clearly $\|A_n(x)\| = \|x\|$ for all $x\in\ell^2$, so that $\|A_n\| = 1$ for all $n\ge 1$. In particular, each $A_n$ is a bounded linear functional. 

We claim that $A_n\to 0$ in the weak operator topology. To this end, we must show that for every $x = (x_n)_{n = 1}^\infty\in\ell^2$ and $\varphi\in\left(\ell^2\right)^\ast$, $\varphi(A_n x)\to 0$. Indeed, due to the Riesz Representation Theorem, there is a $y = (y_n)_{n = 1}^\infty\in\ell^2$ such that $\varphi = \langle\cdot, y\rangle$. We have 
\begin{align*}
    |\varphi(A_n x)| &= \big|\big\langle (\underbrace{0, \dots, 0}_{n\text{ times}}, x_1, x_2, \dots), (y_1, y_2, \dots)\big\rangle\big|\\
    &= \big|\big\langle(0,\dots, 0, x_1, x_2, \dots), (0, \dots, 0, y_{n + 1}, y_{n + 2}, \dots)\big\rangle\big|\\
    &\le\|x\|_2\left(\sum_{k = n + 1}^\infty |y_k|^2\right)^{\frac{1}{2}}.
\end{align*}
Since the sum $\displaystyle\sum_{k = 1}^\infty |y_k|^2$ converges, the tail-sum goes to $0$ as $n\to\infty$, and hence, $|\varphi(A_n x)|\to 0$, i.e., $\varphi(A_n x)\to 0$ as $n\to\infty$. Thus $A_n\to 0$ in the weak operator topology.

On the other hand, we claim that $A_n$ does not converge to $0$ in the strong operator topology. Indeed, consider $x = (1, 0, 0, \dots)$, so that 
\begin{equation*}
    A_n x = (\underbrace{0, \dots, 0}_{n\text{ times}}, 1, 0, 0, \dots),
\end{equation*}
in particular, $A_n x\not\to0$ as $n\to\infty$, since $\|A_n x\| = 1$ for all $n\ge 1$. Hence $A_n$ does not converge to $0$ in the strong operator topology.

\subsection{Part (c)}

Note that $A_n^t\to A^t$ in the weak operator topology if and only if for every $x^\ast\in X^\ast$ and $\varphi\in X^{\ast\ast}$, the sequence $\varphi(A_n^t x^\ast)$ converges to $\varphi(A^t x^\ast)$ in $\K$. Since $X$ is reflexive, the natural map $J\colon X\to X^{\ast\ast}$ is surjective. Therefore, there exists a $y\in X$ such that $\varphi = J(x)$, so that 
\begin{equation*}
    \varphi\left(A_n^t x^\ast\right) = J(y)\left(A_n^t x^\ast\right) = \left(A_n^t x^\ast\right)(y) = x^\ast\left(A_n y\right),
\end{equation*}
and similarly, $\varphi\left(A^t x^\ast\right) = x^\ast(A y)$.

Since $A_n\to A$ in the weak operator topology of $\calL(X)$, the sequence $x^\ast(A_n y)$ converges to $x^\ast(Ay)$ in $\K$ because $x^\ast\in X^\ast$ and $y\in X$. This shows that $A_n^t\to A^t$ in the weak operator topology of $\calL(X^\ast)$.

\subsection{Part (d)}

First we show that $A_n\to 0$ in the strong operator topology. Let $x\in\ell^2$. We have seen in Problem 2 (b) that $e_n\to 0$ in the weak topology of $\ell^2$. Since $\langle \cdot, x\rangle$ is a bounded linear functional on $\ell^2$, this means, in particular, that $\langle e_n, x\rangle\to 0$ as $n\to\infty$. Therefore, $\langle x, e_n\rangle = \overline{\langle e_n, x\rangle}\to 0$ as $n\to\infty$. Hence, $A_nx = \langle x, e_n\rangle e_1\to 0$ as $n\to\infty$, i.e., $A_n\to 0$ in the strong operator topology.

Next, suppose for the sake of contradiction that $A_n^t\to 0$ in $\calL((\ell^2)^\ast)$ in the strong operator topology. Then, for every $x^\ast\in\left(\ell^2\right)^\ast$, we must have $A_n^t x^\ast\to 0$ in $\left(\ell^2\right)^\ast$. Due to the Riesz Representation Theorem, every linear functional on $\ell^2$ is of the form $\langle\cdot, x\rangle$, so that $x^\ast = \langle\cdot, x\rangle$ for some $x\in \ell^2$. Therefore, $A_n^t x^\ast = \langle A_n\cdot, x\rangle = \langle \cdot, A_n^\ast x\rangle$. In particular, this means $\|A_n^t x^\ast\| = \|\langle\cdot, A_n^\ast x\rangle\| = \|A_n^\ast x\|$.

Next, we claim that $A_n^\ast x = \langle x, e_1\rangle e_n$. Indeed, for any $x = (x_n)_{n\ge 1}, y = (y_n)_{n\ge 1}\in\ell^2$, we have 
\begin{equation*}
    \langle A_n x, y\rangle = \langle \langle x, e_n\rangle e_1, y\rangle = \langle x_n e_1, y\rangle = x_n\overline y_1 = \langle x, y_1 e_n\rangle = \langle x, \langle y, e_1\rangle e_n\rangle = \langle x, A_n^\ast y\rangle,
\end{equation*}
as desired. Therefore, $\|A_n^\ast x\| = \|\langle x, e_1\rangle e_n\| = |\langle x, e_1\rangle|$.

Choosing $x = e_1$, that is, choosing $x^\ast = \langle\cdot, e_1\rangle$, we note that $\|A_n^t x^\ast\| = |\langle e_1, e_1\rangle| = 1$ for all $n\ge1$, and hence, $A_n^t x^\ast\not\to 0$ in $\left(\ell^2\right)^\ast$. Hence, $A_n^t\not\to0$ in the strong operator topology, as desired.
\end{document}